\section{A common BFGS energy principle}\label{art:energy}

The useful comparison quantity measures displacement in the current inverse
metric. Its change separates progress toward a reference point from the cost of
the new secant. We first derive this identity, and then show that the BFGS
update forces enough reciprocal growth to rule out a persistent objective gap
when the secant cost is controlled.

\begin{definition}[Admissible runs]\label{art:admissible}
Let $F:\R^n\to\R$ be finite and convex, with $n\ge1$.
An admissible run is an infinite sequence of differentiable accepted points
$x_k$, nonzero gradients $g_k=\nabla F(x_k)$, and matrices generated by
\begin{align}
 s_k&=-\alpha_kH_kg_k,\qquad x_{k+1}=x_k+s_k,
       \qquad y_k=g_{k+1}-g_k,\notag\\
 H_{k+1}&=\left(I-\frac{s_ky_k^T}{t_k}\right)
 H_k\left(I-\frac{y_ks_k^T}{t_k}\right)
       +\frac{s_ks_k^T}{t_k},\qquad t_k=y_k^Ts_k,\label{art:update}
\end{align}
where $H_0\succ0$, $\alpha_k>0$, and fixed constants $0<c_1<c_2<1$ satisfy
\begin{equation}\label{art:aw}
 F(x_{k+1})\le F(x_k)+c_1g_k^Ts_k,
 \qquad g_{k+1}^Ts_k\ge c_2g_k^Ts_k.
\end{equation}
\end{definition}
These are the usual actual-gradient, full-BFGS conventions
\cite{LO}; see also \preprintcite{XW}.
A zero gradient is successful termination. Statements about
admissible runs do not assert that a differentiable acceptable point exists on
every ray, or that every implementation finds one.

Write
\[
 \begin{gathered}
 B_k=H_k^{-1},\qquad q_k=g_k^TH_kg_k,\qquad
 d_k=-g_k^Ts_k=\alpha_kq_k,\\
 \Delta_k=F(x_k)-F(x_{k+1}),\qquad c=1-c_2.
 \end{gathered}
\]
\begin{lemma}\label{art:run-invariants}
Every admissible run satisfies $H_k\succ0$, $q_k,d_k,t_k>0$ and
\begin{equation}\label{art:decrease-bounds}
 c_1d_k\le\Delta_k\le d_k,\qquad t_k\ge cd_k.
\end{equation}
Moreover,
\begin{equation}\label{art:hessian-update}
 B_{k+1}=B_k-\frac{g_kg_k^T}{q_k}+\frac{y_ky_k^T}{t_k},
 \qquad M_k:=B_k-\frac{g_kg_k^T}{q_k}\succeq0,
 \qquad M_ks_k=0.
\end{equation}
If the objective values have a finite limit, then $\sum_kd_k<\infty$.
\end{lemma}
\begin{proof}
Given $H_k\succ0$, nonzero $g_k$ gives $q_k,d_k>0$; weak Wolfe gives
$t_k\ge(1-c_2)d_k>0$. For any vector $v$, the quadratic form in
\eqref{art:update} is a sum of two nonnegative terms. If both vanish, then
$s_k^Tv=0$ and $(I-y_ks_k^T/t_k)v=0$, whence $v=0$.
Thus $H_{k+1}\succ0$, proving the induction. Armijo gives the first lower
bound in \eqref{art:decrease-bounds}; the supporting inequality for $F$ gives
the upper bound. Summing the lower bound proves summability when the values
have a finite limit.

The Hessian identity is the classical equivalent form of BFGS
\cite[Section~5]{FND_N92}; it also follows by multiplying its right side by
\eqref{art:update}, using $B_ks_k=-\alpha_kg_k$ and $s_k^Ty_k=t_k$.
Finally, metric Cauchy--Schwarz gives
$(g_k^Tv)^2\le q_kv^TB_kv$, and
$g_k^Ts_k=-\alpha_kq_k$ gives $M_ks_k=0$.
\end{proof}

\begin{lemma}[Reference energy]\label{art:reference-energy}
For any fixed $x_*$, put $r_k=x_k-x_*$ and
\[
 E_k=r_k^TB_kr_k,\qquad a_k=g_k^Tr_k,\qquad
 R_k=\frac{(y_k^Tr_{k+1})^2}{t_k}.
\]
Then
\begin{equation}\label{art:energy-identity}
 E_{k+1}=E_k-\frac{a_k^2}{q_k}+R_k.
\end{equation}
If $x_*$ minimizes $F$, with $F_*=F(x_*)$, then
\begin{equation}\label{art:energy-budget}
 (F(x_N)-F_*)^2\sum_{k<N}q_k^{-1}
       \le E_0+\sum_{k<N}R_k.
\end{equation}
\end{lemma}
\begin{proof}
Since $r_{k+1}=r_k+s_k$ and $M_ks_k=0$, the Hessian update gives
\[
 r_{k+1}^TB_{k+1}r_{k+1}
 =r_k^TM_kr_k+\frac{(y_k^Tr_{k+1})^2}{t_k},
\]
which is \eqref{art:energy-identity}. At a minimizer, convexity gives
$a_k\ge F(x_k)-F_*$; the values decrease. Summation and $E_N\ge0$ yield
\eqref{art:energy-budget}.
\end{proof}

The negative term is progress measured through a supporting gradient. The
positive term records the secant's component in the next reference direction.
An upper bound on this correction becomes useful only when the reciprocal
sum is known to grow sufficiently fast. That growth is forced by the full
matrix update, without a spectral bound on the matrices.

\subsection{Actual principal minors and reciprocal growth}

For a symmetric $n\times n$ matrix $H$, define
\[
 e_j(H)=\sum_{\substack{I\subset\{1,\ldots,n\}\\|I|=j}}\det H_{I,I},
 \qquad e_{n+1}(H)=0.
\]
These are actual principal determinants. Multilinearity of the determinant
identifies them with the coefficients of $\det(I+zH)$; hence they are invariant
under orthogonal changes of coordinates and equal the elementary symmetric
polynomials of the eigenvalues. In particular $e_1=\tr$ and $e_n=\det$.
The determinant endpoint is classical; trace/determinant potentials have long
been used in BFGS analysis \cite{FND_N92,FND_BN89}. We need all adjacent degrees.

\begin{theorem}[Principal-minor hierarchy]\label{art:minor-hierarchy}
Let $H\succ0$, $s=-\alpha Hg$, $\alpha>0$, and suppose
\[
 q=g^THg>0,\qquad \beta=-\frac{y^THg}{q}>0.
\]
For the inverse-BFGS update $H_+$ in \eqref{art:update},
\begin{equation}\label{art:hierarchy}
 e_j(H_+)\le\frac{\norm{y}^2}{\beta^2q}e_{j+1}(H)
                   +\frac\alpha\beta e_j(H),\qquad 1\le j\le n.
\end{equation}
At $j=n$ the exact identity is $\det H_+=(\alpha/\beta)\det H$.
\end{theorem}
\begin{proof}
Let $Z$ have $j$ orthonormal columns. The complete compression calculation in
Lemma~\ref{art:gram-compression} gives
\[
 \det(Z^TH_+Z)\le
 \frac{\det\operatorname{Gram}_H(y,Z)}{\beta^2q}
       +\frac\alpha\beta\det(Z^THZ),
\]
where $\operatorname{Gram}_H(y,Z)=[y\ Z]^TH[y\ Z]$.
Choose an orthonormal eigenbasis of $H$, with eigenvalues $\lambda_i>0$, and
write $y_i$ for the coordinates of $y$ in this basis. Summing over its
$j$-column selections $Z_I$ gives
\begin{align*}
 \sum_{|I|=j}\det\operatorname{Gram}_H(y,Z_I)
 &=\sum_{|I|=j}\left(\prod_{i\in I}\lambda_i\right)
                  \sum_{\ell\notin I}\lambda_\ell y_\ell^2\\
 &=\sum_{|J|=j+1}\left(\prod_{i\in J}\lambda_i\right)
                  \sum_{\ell\in J}y_\ell^2
 \le\norm{y}^2e_{j+1}(H).
\end{align*}
The first equality is the Schur complement of the diagonal compression;
the second groups terms by $J=I\cup\{\ell\}$. Orthogonal invariance now
proves \eqref{art:hierarchy}. The last assertion follows from the full-size
case in Lemma~\ref{art:gram-compression}.
\end{proof}

\begin{theorem}[Reciprocal growth]\label{art:reciprocal-growth}
If an admissible run has bounded gradients and $d_k\to0$, then
\begin{equation}\label{art:reciprocal-limits}
 \left(\prod_{k<N}q_k^{-1}\right)^{1/N}\longrightarrow+\infty,
 \qquad \frac1N\sum_{k<N}q_k^{-1}\longrightarrow+\infty.
\end{equation}
\end{theorem}
\begin{proof}
Choose $G>0$ with $\norm{g_k}\le G$, and put $K_0=4G^2/c^2$.
For the actual run,
\[
 \beta_k=\frac{t_k}{d_k}\ge c,\qquad
 \norm{y_k}\le2G,\qquad \alpha_k=\frac{d_k}{q_k}.
\]
Theorem~\ref{art:minor-hierarchy} therefore yields the normalized hierarchy
\begin{equation}\label{art:normalized-hierarchy}
 q_ke_j(H_{k+1})\le
 K_0e_{j+1}(H_k)+\frac{d_k}{c}e_j(H_k).
\end{equation}
This recurrence is derived for the run's principal determinants.

Fix $\eta>0$ and choose $J$ so that $d_k/c\le\eta$ for $k\ge J$.
For $H\succ0$ define
\[
 V_\eta(H)=\sum_{j=1}^n(K_0/\eta)^{j-1}e_j(H)>0.
\]
Multiply \eqref{art:normalized-hierarchy} by these weights and sum. The
$e_{j+1}$ terms contribute at most $\eta V_\eta(H_k)$, as do the remaining
terms. Also $q_k\le G^2\tr H_k\le G^2V_\eta(H_k)$. Consequently
\[
 V_\eta(H_{k+1})\le\frac{2\eta}{q_k}V_\eta(H_k),\qquad
 \prod_{k=J}^{J+m-1}q_k^{-1}
 \ge\frac{1}{G^2V_\eta(H_J)(2\eta)^{m-1}}.
\]
For the last bound, iterate through $m-1$ updates and apply the bound on
$q_{J+m-1}$ at the remaining endpoint. Taking $m$th roots gives a lower
limit of at least $(2\eta)^{-1}$. A fixed prefix contributes a factor tending
to one. Since $\eta$ is arbitrary, the geometric mean diverges; AM--GM gives
the arithmetic-mean conclusion.
\end{proof}

Product estimates and AM--GM are established tools in BFGS analysis
\cite[Theorem~3.7]{S16}\cite[Appendices C--D]{FND_JJM24AW}.
Here the normalized hierarchy supplies them in the bounded-gradient,
vanishing-decrease setting. Neither limit in \eqref{art:reciprocal-limits}
asserts that $q_k\to0$ pointwise or that $H_k$ converges.

\subsection{Moving comparison directions}

A reference direction can be allowed to move, provided that it continues to
see descent, controls the incoming secant, and does not increase its energy
too rapidly. The identity $M_ks_k=0$ shows that movement along the executed
step has no cost in the part of the update that remains after the negative
rank-one term. A positive-semidefinite comparison field expresses these
requirements without a choice of frame.

For differentiable $x,z$, write $g=\nabla F(x)$, $h=\nabla F(z)$,
$s=z-x$, $y=h-g$, $d=-g^Ts$, $\Delta=F(x)-F(z)$ and $t=y^Ts$.
Call the pair eligible when
\[
 d>0,\qquad c_1d\le\Delta\le d,\qquad t\ge cd.
\]
All accepted pairs are eligible. Requiring a property of every eligible pair
is an objective assumption; it does not assert that an arbitrary sequence of
eligible pairs can be realized by BFGS. Put
$\Pi_s=I-ss^T/(s^Ts)$.

\begin{definition}[Band certificate]\label{art:band-certificate}
Let $K=\{F\le b\}$ be compact, with $b>F_*=\min F$.
For each nonempty $K_\eta=\{x\in K:F(x)\ge F_*+\eta\}$, $\eta>0$,
a band certificate consists of a finite-valued symmetric PSD field
$W_\eta(x)$ on its differentiable points and constants
$m_\eta>0$, $C_\eta\ge0$, $L_\eta\ge0$, $\delta_\eta>0$ such that
\begin{align}
 g^TW_\eta(x)g&\ge m_\eta,\label{art:visibility}\\
 y^TW_\eta(z)y&\le C_\eta t,\label{art:secant-control}\\
 \Pi_sW_\eta(z)\Pi_s&\preceq
 (1+L_\eta d)\Pi_sW_\eta(x)\Pi_s.\label{art:transport}
\end{align}
The first inequality holds at every differentiable point of $K_\eta$;
the last two hold for all eligible pairs in $K_\eta$ with $d\le\delta_\eta$.
\end{definition}
Neither continuity, uniform boundedness nor positive definiteness of $W_\eta$
is required. The constants and field may depend on the positive band.
Allowing an arbitrary real secant constant gives the same condition, since
it can always be replaced by its maximum with zero.

\begin{theorem}[Transported comparison certificate]\label{art:certificate}
If $K$ admits the band certificates above, every admissible run with $x_0\in K$
satisfies
\[
 F(x_k)\to F_*,\qquad
 \operatorname{dist}(x_k,\arg\min F)\to0.
\]
If the minimizer is unique, the iterates converge to it.
\end{theorem}
\begin{proof}
Armijo keeps the run in $K$, and its values have a finite limit.
Lemma~\ref{art:run-invariants} gives $\sum d_k<\infty$.
Gradients are bounded there: choose $\epsilon>0$ and bound the continuous
finite convex function on $K+\epsilon\overline B$. Its supporting inequality
at $x$ evaluated at $x+\epsilon u$, for every unit $u$, bounds every
subgradient at $x$. Thus Theorem~\ref{art:reciprocal-growth} applies.

Suppose the limiting gap is $\ell>0$. On a tail $k\ge J$, use the field
on $K_{\ell/2}$ and ensure $d_k\le\delta_{\ell/2}$; omit the band subscripts.
Set $E_k=\tr(B_kW(x_k))\ge0$. Since
$M_k=\Pi_{s_k}M_k\Pi_{s_k}\succeq0$, trace monotonicity and
\eqref{art:secant-control}--\eqref{art:transport} give
\begin{align*}
 E_{k+1}
 &=\tr(M_kW(x_{k+1}))+
       \frac{y_k^TW(x_{k+1})y_k}{t_k}\\
 &\le(1+Ld_k)\tr(M_kW(x_k))+C
 \le(1+Ld_k)(E_k-m/q_k)+C.
\end{align*}
Trace monotonicity here follows from
$\tr(AQ)=\tr(Q^{1/2}AQ^{1/2})\ge0$ for PSD $A,Q$.
Let $P_J=1$ and $P_{k+1}=P_k(1+Ld_k)$. Summability gives
$1\le P_k\le\exp(L\sum_{j\ge J}d_j)=:\overline P$.
Dividing by $P_{k+1}$, summing and using $E_N\ge0$ yields
\[
 \frac{m}{\overline P}\sum_{k=J}^{N-1}q_k^{-1}
       \le E_J+C(N-J),
\]
contradicting \eqref{art:reciprocal-limits}. This normalization is the usual
summable multiplicative-error argument; compare \cite[Lemmas~2.2--2.3]{FND_CV13}.
Finally, a sequence in compact $K$ with values tending to $F_*$ cannot stay
a positive distance from the closed solution set: a convergent subsequence
would have a nonoptimal limit with value $F_*$. With a unique minimizer the
same argument gives convergence of the entire sequence.
\end{proof}
